\appendix
\setcounter{table}{0}
\renewcommand{\thetable}{A.\arabic{table}}
\renewcommand{\theHtable}{A.\arabic{table}}

\section{Definition of Evaluation Metrics}
\label{app:evaluation_metrics}

This appendix defines the evaluation metrics that are specific to surrogate-assisted candidate screening, probability calibration, ranking-based reinforcement prediction, and optimization efficiency. Standard classification and regression metrics, such as ROC-AUC, PR-AUC, F1 score, MAE, RMSE, $R^2$, and MAPE, follow their conventional definitions and are not repeated here.

\subsection{Screening Metrics for Feasibility Prediction}
\label{app:screening_metrics}

In this study, the feasibility surrogate is used as a conservative screening model. Therefore, its performance is evaluated not only by ordinary classification metrics, but also by whether it can reject low-value candidates while retaining most truly feasible candidates.

Let $y_i \in {0,1}$ denote the true feasibility label of candidate $i$, where $y_i=1$ indicates a feasible candidate. Let $\hat{p}_{f,i}$ denote the predicted or calibrated feasibility probability, and let $\tau$ denote the screening threshold. A candidate is retained for potential FEA evaluation if
\begin{equation}
\hat{p}_{f,i} \geq \tau .
\end{equation}

The feasible recall is defined as
\begin{equation}
R_{\mathrm{fea}}
=
\frac{
\sum_{i=1}^{N}
\mathbb{I}(y_i=1)\mathbb{I}(\hat{p}_{f,i}\geq \tau)
}{
\sum_{i=1}^{N}
\mathbb{I}(y_i=1)
},
\label{eq:app_feasible_recall}
\end{equation}
where $\mathbb{I}(\cdot)$ is the indicator function. This metric measures the proportion of truly feasible candidates retained after surrogate screening. A high feasible recall is required because falsely filtering out feasible candidates may reduce the chance of finding a valid optimal design.

The reject rate is defined as
\begin{equation}
R_{\mathrm{rej}}
=
\frac{
\sum_{i=1}^{N}
\mathbb{I}(\hat{p}_{f,i}<\tau)
}{N}.
\label{eq:app_reject_rate}
\end{equation}
It measures the proportion of candidates rejected before true FEA evaluation. A higher reject rate indicates greater potential for reducing FEA calls, but it is meaningful only when feasible recall remains sufficiently high.

The false reject count is defined as
\begin{equation}
N_{\mathrm{fr}}
=
\sum_{i=1}^{N}
\mathbb{I}(y_i=1)\mathbb{I}(\hat{p}_{f,i}<\tau).
\label{eq:app_false_reject}
\end{equation}
It directly counts the number of truly feasible candidates incorrectly rejected by the screening rule. This metric is reported because even a small false reject rate may be important when the number of feasible candidates is limited.

\subsection{Probability Calibration Metrics}
\label{app:calibration_metrics}

The feasibility surrogate outputs a probability rather than a binary decision. In surrogate-assisted optimization, probability reliability is important because the screening threshold is selected based on predicted feasibility probabilities.

The Brier score is used to measure the mean squared error between predicted probabilities and true binary labels:
\begin{equation}
\mathrm{Brier}
=
\frac{1}{N}
\sum_{i=1}^{N}
(\hat{p}_{f,i}-y_i)^2.
\label{eq:app_brier}
\end{equation}
A lower Brier score indicates better probabilistic prediction.

Expected calibration error (ECE) measures the discrepancy between predicted probability and empirical feasibility frequency. The predictions are divided into $M$ probability bins. Let $B_m$ denote the set of samples in bin $m$. The average confidence and empirical feasible rate in bin $m$ are defined as
\begin{equation}
\mathrm{conf}(B_m)
=
\frac{1}{|B_m|}
\sum_{i\in B_m}
\hat{p}_{f,i},
\label{eq:app_confidence}
\end{equation}
\begin{equation}
\mathrm{freq}(B_m)
=
\frac{1}{|B_m|}
\sum_{i\in B_m}
y_i.
\label{eq:app_frequency}
\end{equation}
The ECE is then calculated as
\begin{equation}
\mathrm{ECE}
=
\sum_{m=1}^{M}
\frac{|B_m|}{N}
\left|
\mathrm{freq}(B_m)
-
\mathrm{conf}(B_m)
\right|.
\label{eq:app_ece}
\end{equation}
A smaller ECE indicates that the predicted feasibility probabilities are more consistent with the observed feasible frequencies.

\subsection{Ranking Metrics for Reinforcement Prediction}
\label{app:ranking_metrics}

The reinforcement surrogate is mainly used to rank candidates by material economy before true FEA evaluation. Therefore, ranking-oriented metrics are reported in addition to regression errors.

Pairwise ranking accuracy evaluates whether the surrogate correctly predicts the relative order of reinforcement usage between candidate pairs. Given two candidates $i$ and $j$, the ordering is correctly predicted if
\begin{equation}
(m_{s,i}-m_{s,j})(\hat{m}_{s,i}-\hat{m}_{s,j}) > 0 .
\end{equation}
The pairwise ranking accuracy is defined as
\begin{equation}
A_{\mathrm{rank}}
=
\frac{
\sum_{i<j}
\mathbb{I}
\left[
(m_{s,i}-m_{s,j})(\hat{m}_{s,i}-\hat{m}_{s,j}) > 0
\right]
}{
\binom{N}{2}
}.
\label{eq:app_pairwise_ranking}
\end{equation}
A higher value indicates that the surrogate better preserves the relative material economy of candidates.

Spearman's rank correlation coefficient is also used to measure the monotonic agreement between true and predicted reinforcement usage:
\begin{equation}
\rho_s
=
1
-
\frac{
6\sum_{i=1}^{N} d_i^2
}{
N(N^2-1)
},
\label{eq:app_spearman}
\end{equation}
where $d_i$ is the difference between the true rank and predicted rank of candidate $i$. A value closer to 1 indicates better rank consistency.

Top-$k$ recall evaluates whether the surrogate can identify the truly low-reinforcement candidates. Let $\mathcal{T}_k$ denote the set of candidates with the lowest true reinforcement usage, and let $\hat{\mathcal{T}}_k$ denote the set of candidates with the lowest predicted reinforcement usage. The top-$k$ recall is defined as
\begin{equation}
R_{\mathrm{top}\text{-}k}
=
\frac{
|\mathcal{T}_k \cap \hat{\mathcal{T}}_k|
}{
|\mathcal{T}_k|
}.
\label{eq:app_topk_recall}
\end{equation}
This metric is particularly relevant for optimization because only a small subset of promising candidates is selected for true FEA evaluation.

The prediction bias of reinforcement usage is defined as
\begin{equation}
\mathrm{Bias}
=
\frac{1}{N}
\sum_{i=1}^{N}
(\hat{m}_{s,i}-m_{s,i}).
\label{eq:app_bias}
\end{equation}
A positive value indicates overestimation, while a negative value indicates underestimation. Bias is reported because systematic overestimation or underestimation may affect candidate ranking and FEA budget allocation.

\subsection{Optimization Efficiency and Solution Quality Metrics}
\label{app:optimization_metrics}

In the optimization experiments, the primary computational metric is the number of true FEA calls:
\begin{equation}
N_{\mathrm{FEA}}
=
\sum_{t}
|\mathcal{S}_t|,
\label{eq:app_fea_calls}
\end{equation}
where $\mathcal{S}_t$ is the set of candidates submitted to the true evaluator at iteration $t$.

The FEA reduction ratio is calculated relative to full FEA optimization:
\begin{equation}
R_{\mathrm{FEA}}
=
1
-
\frac{
N_{\mathrm{FEA}}^{\mathrm{sur}}
}{
N_{\mathrm{FEA}}^{\mathrm{full}}
},
\label{eq:app_fea_reduction}
\end{equation}
where $N_{\mathrm{FEA}}^{\mathrm{sur}}$ and $N_{\mathrm{FEA}}^{\mathrm{full}}$ are the numbers of true FEA calls used by the surrogate-assisted strategy and the full FEA baseline, respectively.

The feasible-solution discovery rate is defined as the proportion of optimization runs that successfully find at least one feasible design:
\begin{equation}
R_{\mathrm{succ}}
=
\frac{
N_{\mathrm{runs}}^{\mathrm{feasible}}
}{
N_{\mathrm{runs}}
}.
\label{eq:app_success_rate}
\end{equation}
This metric reflects whether surrogate screening reduces the chance of discovering feasible solutions.

For runs that find feasible solutions, the final material cost of the best FEA-verified feasible design is reported. To compare surrogate-assisted optimization with full FEA optimization under the same layout and random seed, the normalized objective ratio is calculated as
\begin{equation}
R_{\mathrm{obj}}
=
\frac{
C_{\mathrm{best}}^{\mathrm{sur}}
}{
C_{\mathrm{best}}^{\mathrm{full}}
},
\label{eq:app_objective_ratio}
\end{equation}
where $C_{\mathrm{best}}^{\mathrm{sur}}$ and $C_{\mathrm{best}}^{\mathrm{full}}$ are the best verified feasible costs obtained by the surrogate-assisted method and the full FEA baseline, respectively. A value close to 1 indicates that the surrogate-assisted method achieves comparable solution quality.

The number of FEA calls required to find the first feasible solution is also reported:
\begin{equation}
N_{\mathrm{first}}
=
\min
\left\{
n
\mid
\exists i \leq n,\ y_{f,i}=1
\right\}.
\label{eq:app_first_feasible}
\end{equation}
This metric indicates how quickly an optimization strategy reaches the feasible region under a limited FEA budget.


\section{Additional Calibration Results}
\label{app:calibration_results}
\setcounter{table}{0}
\renewcommand{\thetable}{B.\arabic{table}}
\renewcommand{\theHtable}{B.\arabic{table}}

Tables~\ref{tab:appendix_feasibility_calibrator} and~\ref{tab:appendix_steel_calibrator} report the calibration-model selection results for feasibility probability calibration and steel-usage residual calibration under different local sample sizes. For each sample size, several lightweight calibrators are trained and the model with the best Brier score is reported together with other evaluation metrics. The results show that the most suitable calibrator varies with the amount of available local calibration data.

\begin{table}[H]
\centering
\caption{Best feasibility probability calibrator under different local sample sizes.}
\label{tab:appendix_feasibility_calibrator}
\small
\setlength{\tabcolsep}{6pt}
\begin{tabularx}{\textwidth}{@{}>{\centering\arraybackslash}p{0.12\textwidth}
>{\raggedright\arraybackslash}X
>{\centering\arraybackslash}p{0.10\textwidth}
>{\centering\arraybackslash}p{0.09\textwidth}
>{\centering\arraybackslash}p{0.13\textwidth}
>{\centering\arraybackslash}p{0.13\textwidth}@{}}
\toprule
Calib. samples & Best feasibility calibrator & Brier & ECE & Screen reject & Screen recall \\
\midrule
25 & Platt scaling & 0.0505 & 0.0455 & 47.4\% & 0.9476 \\
50 & Residual random forest & 0.0464 & 0.0406 & 51.6\% & 0.9485 \\
100 & Residual random forest & 0.0420 & 0.0325 & 52.3\% & 0.9487 \\
200 & Residual random forest & 0.0381 & 0.0261 & 59.1\% & 0.9481 \\
500 & Residual random forest & 0.0344 & 0.0201 & 63.9\% & 0.9441 \\
\bottomrule
\end{tabularx}
\end{table}

\begin{table}[H]
\centering
\caption{Best steel residual calibrator under different local sample sizes.}
\label{tab:appendix_steel_calibrator}
\small
\setlength{\tabcolsep}{6pt}
\begin{tabularx}{\textwidth}{@{}>{\centering\arraybackslash}p{0.12\textwidth}
>{\raggedright\arraybackslash}X
>{\centering\arraybackslash}p{0.12\textwidth}
>{\centering\arraybackslash}p{0.12\textwidth}
>{\centering\arraybackslash}p{0.08\textwidth}
>{\centering\arraybackslash}p{0.11\textwidth}@{}}
\toprule
Calib. samples & Best steel residual model & MAE (kg) & RMSE (kg) & $R^2$ & Pair acc. \\
\midrule
25 & Ridge & 2395.9 & 3162.8 & 0.9924 & 0.9769 \\
50 & Gaussian process & 2229.1 & 3003.6 & 0.9932 & 0.9785 \\
100 & Gaussian process & 2044.7 & 2756.0 & 0.9943 & 0.9797 \\
200 & LightGBM & 1741.3 & 2386.0 & 0.9957 & 0.9821 \\
500 & LightGBM & 1444.8 & 1992.6 & 0.9970 & 0.9851 \\
\bottomrule
\end{tabularx}
\end{table}
